% test-016.tex -- comandos \sprad, \spdiam, \spcirc, \spAcirc, \spPcirc y \spVol
%
% Compilar:  lualatex-dev test-016.tex
% Validar:   show-pdf-tags --xml test-016.pdf | rnv-wrapp
%            verapdf --flavour ua2 --format text test-016.pdf
%
% Cada caso es un parrafo numerado, para poder referirse a el al escuchar
% el PDF. Los comentarios "doc:" indican lo que documenta el manual; en el
% resto, la lectura se revisa escuchando. Todos trabajan en modo
% matematico. \sprad, \spdiam y \spVol no llevan argumento, solo llaves.
% \spcirc recibe, entre parentesis, un centro o un centro y un radio (un
% entero o decimal, una letra, una medida con unidad o dos puntos). \spAcirc y
% \spPcirc reciben un centro o un numero, nunca una coma.
%
% Entradas que no van aqui porque detienen la compilacion: una expresion
% como radio (\sqrt{2}, r+1, \pi) y una coma en \spAcirc o \spPcirc. Un
% radio decimal (2.5 y 2,5) requiere spcirc-decimal.patch.
%
% Los casos con potencias (\sprad^{2}) van al final, en su propia seccion,
% porque combinan un comando de este paquete con un exponente.
\DocumentMetadata
  {
    lang=es-CL,
    pdfversion=2.0,
    pdfstandard=ua-2,
    tagging=on,
    tagging-setup={math/setup={mathml-SE}},
  }
\documentclass{article}
\usepackage[spanish]{babel}
\usepackage{unicode-math}
\usepackage{spintent}
\usepackage{hyperref}
\hypersetup
  {
    pdfdisplaydoctitle = true,
    pdftitle           = {test-016: sprad, spdiam, spcirc, spAcirc, spPcirc y spVol},
  }
\pagestyle{empty}
\begin{document}

\section*{Comando sprad}

% doc: «radio»
1. Sin llaves: $\sprad$.

% doc: «radio»
2. Con concept true, el valor por defecto: $\sprad[concept=true]$.

% doc: «r»
3. Con concept false: $\sprad[concept=false]$.

% doc: «r»
4. Con concept false y read-cmd school: $\sprad[concept=false,read-cmd=school]$.

% doc: «r»
5. Con concept false y read-cmd mathcat: $\sprad[concept=false,read-cmd=mathcat]$.

% doc: «radio»
6. Con upper-case false, el valor por defecto: $\sprad[upper-case=false]$.

% doc: «radio»
7. Con upper-case true: $\sprad[upper-case=true]$.

% doc: «R»
8. Con upper-case true y concept false: $\sprad[upper-case=true,concept=false]$.

% doc: «r»
9. Con upper-case true, concept false y read-cmd mathcat: $\sprad[upper-case=true,concept=false,read-cmd=mathcat]$.

% doc: «radio»
10. Con print-sym false, el valor por defecto: $\sprad[print-sym=false]$.

% doc: «radio»
11. Con print-sym true: $\sprad[print-sym=true]$.

% doc: «radio 1»
12. Con sub-index: $\sprad[sub-index={1}]$.

% doc: «radio 12»
13. Con sub-index de dos cifras: $\sprad[sub-index={12}]$.

% doc: «radio sub 1»
14. Con sub-index y read-sub true: $\sprad[sub-index={1},read-sub=true]$.

% doc: «radio 1»
15. Con sub-index y read-sub false: $\sprad[sub-index={1},read-sub=false]$.

% doc: «r 1»
16. Con sub-index y concept false: $\sprad[sub-index={1},concept=false]$.

% doc: «radio medida exacta»
17. Con read-arg: $\sprad[read-arg={medida exacta}]$.

% doc: «medida exacta»
18. Con read-arg y concept false: $\sprad[concept=false,read-arg={medida exacta}]$.

% doc: «radio medida exacta»
19. Con read-arg y sub-index: $\sprad[read-arg={medida exacta},sub-index={1},read-sub=true]$.

% doc: «radio»
20. Con silent false, el valor por defecto: $\sprad[silent=false]$.

% doc: «»
21. Con silent true: $\sprad[silent=true]$.

% doc: «»
22. Con silent true y sub-index: $\sprad[silent=true,sub-index={1}]$.

% doc: «radio subíndice 1»
23. Con todas las llaves: $\sprad[concept=true,read-cmd=mathcat,upper-case=true,sub-index={1},read-sub=true,silent=false]$.

\section*{Comando spdiam}

% doc: «diámetro»
24. Sin llaves: $\spdiam$.

% doc: «diámetro»
25. Con concept true, el valor por defecto: $\spdiam[concept=true]$.

% doc: «d»
26. Con concept false: $\spdiam[concept=false]$.

% doc: «d»
27. Con concept false y read-cmd school: $\spdiam[concept=false,read-cmd=school]$.

% doc: «d»
28. Con concept false y read-cmd mathcat: $\spdiam[concept=false,read-cmd=mathcat]$.

% doc: «diámetro»
29. Con upper-case false, el valor por defecto: $\spdiam[upper-case=false]$.

% doc: «diámetro»
30. Con upper-case true: $\spdiam[upper-case=true]$.

% doc: «D»
31. Con upper-case true y concept false: $\spdiam[upper-case=true,concept=false]$.

% doc: «d»
32. Con upper-case true, concept false y read-cmd mathcat: $\spdiam[upper-case=true,concept=false,read-cmd=mathcat]$.

% doc: «diámetro»
33. Con print-sym false, el valor por defecto: $\spdiam[print-sym=false]$.

% doc: «diámetro»
34. Con print-sym true: $\spdiam[print-sym=true]$.

% doc: «diámetro 1»
35. Con sub-index: $\spdiam[sub-index={1}]$.

% doc: «diámetro 12»
36. Con sub-index de dos cifras: $\spdiam[sub-index={12}]$.

% doc: «diámetro sub 1»
37. Con sub-index y read-sub true: $\spdiam[sub-index={1},read-sub=true]$.

% doc: «diámetro 1»
38. Con sub-index y read-sub false: $\spdiam[sub-index={1},read-sub=false]$.

% doc: «d 1»
39. Con sub-index y concept false: $\spdiam[sub-index={1},concept=false]$.

% doc: «diámetro medida exacta»
40. Con read-arg: $\spdiam[read-arg={medida exacta}]$.

% doc: «medida exacta»
41. Con read-arg y concept false: $\spdiam[concept=false,read-arg={medida exacta}]$.

% doc: «diámetro medida exacta»
42. Con read-arg y sub-index: $\spdiam[read-arg={medida exacta},sub-index={1},read-sub=true]$.

% doc: «diámetro»
43. Con silent false, el valor por defecto: $\spdiam[silent=false]$.

% doc: «»
44. Con silent true: $\spdiam[silent=true]$.

% doc: «»
45. Con silent true y sub-index: $\spdiam[silent=true,sub-index={1}]$.

% doc: «diámetro subíndice 1»
46. Con todas las llaves: $\spdiam[concept=true,read-cmd=mathcat,upper-case=true,print-sym=true,sub-index={1},read-sub=true,silent=false]$.

\section*{Comando spcirc: argumentos}

% doc: «circunferencia»
47. Sin argumentos: $\spcirc$.

% doc: «circunferencia de centro en O»
48. Con un centro: $\spcirc(O)$.

% doc: «circunferencia de centro en O y radio 3»
49. Con centro y radio entero: $\spcirc(O,3)$.

% doc: «circunferencia de centro en O y radio 3»
50. Con centro y radio, con espacio: $\spcirc(O, 3)$.

% doc: «circunferencia de centro en O prima»
51. Con un centro con prima: $\spcirc(O')$.

% doc: «circunferencia de centro en O sub uno»
52. Con un centro con subíndice: $\spcirc(O_{1})$.

% doc: «circunferencia de centro en P»
53. Con otra letra como centro: $\spcirc(P)$.

% doc: «circunferencia de centro en O y radio 12»
54. Con un radio de dos cifras: $\spcirc(O,12)$.

% doc: «circunferencia de centro en O y radio 2 punto 5»
55. Con un radio decimal con punto: $\spcirc(O,2.5)$.

% doc: «circunferencia de centro en O y radio 2 coma 5»
56. Con un radio decimal con coma: $\spcirc(O,2,5)$.

% doc: «circunferencia de centro en O y radio r»
57. Con un radio dado por una letra: $\spcirc(O,r)$.

% doc: «circunferencia de centro en O y radio r»
58. Con un radio dado por una mayúscula: $\spcirc(O,R)$.

% doc: «circunferencia de centro en O y radio 3 centímetros»
59. Con un radio con medida: $\spcirc(O,3 cm)$.

% doc: «circunferencia de centro en O y radio 3 centímetros»
60. Con la unidad pegada: $\spcirc(O,3cm)$.

% doc: «circunferencia de centro en O y radio 2 punto 5 centímetros»
61. Con una medida decimal: $\spcirc(O,2.5 cm)$.

% doc: «circunferencia de centro en O y radio 10 metros»
62. Con otra unidad: $\spcirc(O,10 m)$.

% doc: «circunferencia de centro en O y radio medida del lado A B»
63. Con un radio dado por un segmento: $\spcirc(O,AB)$.

% doc: «circunferencia de centro en O y radio medida del lado A prima B prima»
64. Con un segmento con primas: $\spcirc(O,A'B')$.

% doc: «circunferencia de centro en O y radio medida del lado A sub uno B sub dos»
65. Con un segmento con subíndices: $\spcirc(O,A_{1}B_{2})$.

% doc: «circunferencia de centro en O prima y radio 3»
66. Con centro con prima y radio: $\spcirc(O',3)$.

\section*{Comando spcirc: llaves}

% doc: «circunferencia de centro en O»
67. Con read-cmd school, el valor por defecto y un centro: $\spcirc[read-cmd=school](O)$.

% doc: «circunferencia de centro en O y radio medida del lado A B»
68. Con read-cmd school, el valor por defecto y un segmento: $\spcirc[read-cmd=school](O,AB)$.

% doc: «circunferencia de centro en o»
69. Con read-cmd mathcat y un centro: $\spcirc[read-cmd=mathcat](O)$.

% doc: «circunferencia de centro en o y radio medida del lado a b»
70. Con read-cmd mathcat y un segmento: $\spcirc[read-cmd=mathcat](O,AB)$.

% doc: «la circunferencia»
71. Con prefix y sin argumentos: $\spcirc[prefix={la}]$.

% doc: «la circunferencia de centro en O»
72. Con prefix y un centro: $\spcirc[prefix={la}](O)$.

% doc: «la circunferencia de centro en O y radio 3»
73. Con prefix y centro y radio: $\spcirc[prefix={la}](O,3)$.

% doc: «circunferencia de centro en O»
74. Con print-sym odot, el valor por defecto: $\spcirc[print-sym=odot](O)$.

% doc: «circunferencia de centro en O»
75. Con print-sym usual: $\spcirc[print-sym=usual](O)$.

% doc: «circunferencia de centro en O y radio 3»
76. Con print-sym usual y radio: $\spcirc[print-sym=usual](O,3)$.

% doc: «circunferencia»
77. Con print-sym usual y sin argumentos: $\spcirc[print-sym=usual]$.

% doc: «circunferencia de centro en O»
78. Con read-sty circun, el valor por defecto: $\spcirc[read-sty=circun](O)$.

% doc: «círculo de centro en O»
79. Con read-sty circle: $\spcirc[read-sty=circle](O)$.

% doc: «círculo de centro en O y radio 3»
80. Con read-sty circle y radio: $\spcirc[read-sty=circle](O,3)$.

% doc: «círculo»
81. Con read-sty circle y sin argumentos: $\spcirc[read-sty=circle]$.

% doc: «circunferencia de centro en O y radio 3»
82. Con diameter false, el valor por defecto: $\spcirc[diameter=false](O, 3)$.

% doc: «circunferencia de centro en O y diámetro 3»
83. Con diameter true: $\spcirc[diameter=true](O, 3)$.

% doc: «circunferencia de centro en O y diámetro medida del lado A B»
84. Con diameter true y un segmento: $\spcirc[diameter=true](O, AB)$.

% doc: «circunferencia de centro en O y diámetro 3 centímetros»
85. Con diameter true y una medida: $\spcirc[diameter=true](O, 3 cm)$.

% doc: «circunferencia de centro en O y radio medida del lado A B»
86. Con radio-sty mside, el valor por defecto: $\spcirc[radio-sty=mside](O, AB)$.

% doc: «circunferencia de centro en O y radio lado A B»
87. Con radio-sty side: $\spcirc[radio-sty=side](O, AB)$.

% doc: «circunferencia de centro en O y radio lado A prima B prima»
88. Con radio-sty side y primas: $\spcirc[radio-sty=side](O, A'B')$.

% doc: «circunferencia de centro en O y radio segmento A B»
89. Con la meta-llave spside: $\spcirc[radio-sty=side,spside={read-sty=segmento}](O, AB)$.

% doc: «circunferencia de centro en O y radio medida del lado A B»
90. Con la meta-llave spmside: $\spcirc[radio-sty=mside,spmside={print-m=true}](O, AB)$.

% doc: «circunferencia de centro en O y radio medida del trazo A B»
91. Con la meta-llave spmside y read-sty: $\spcirc[spmside={read-sty=trazo}](O, AB)$.

% doc: «circunferencia de centro en O y radio 3»
92. Con concept true, el valor por defecto: $\spcirc[concept=true](O,3)$.

% doc: «»
93. Con concept false y sin argumentos: $\spcirc[concept=false]$.

% doc: «O»
94. Con concept false y un centro: $\spcirc[concept=false](O)$.

% doc: «O 3»
95. Con concept false y centro y radio: $\spcirc[concept=false](O,3)$.

% doc: «O medida del lado A B»
96. Con concept false y centro y segmento: $\spcirc[concept=false](O,AB)$.

% doc: «circunferencia 1 de centro en O y radio medida del lado A B»
97. Con sub-index y un segmento: $\spcirc[sub-index={1},radio-sty=mside](O, AB)$.

% doc: «circunferencia 1 de centro en O»
98. Con sub-index y un centro: $\spcirc[sub-index={1}](O)$.

% doc: «circunferencia 1»
99. Con sub-index y sin argumentos: $\spcirc[sub-index={1}]$.

% doc: «circunferencia sub 1 de centro en O»
100. Con sub-index y read-sub true: $\spcirc[sub-index={1},read-sub=true](O)$.

% doc: «circunferencia 1 de centro en O»
101. Con sub-index y read-sub false: $\spcirc[sub-index={1},read-sub=false](O)$.

% doc: «disco de centro en O»
102. Con read-arg: $\spcirc[read-arg={disco}](O)$.

% doc: «disco de centro en O y radio 3»
103. Con read-arg y radio: $\spcirc[read-arg={disco}](O,3)$.

% doc: «disco»
104. Con read-arg y sin argumentos: $\spcirc[read-arg={disco}]$.

% doc: «disco de centro en O»
105. Con read-arg y concept false: $\spcirc[concept=false,read-arg={disco}](O)$.

% doc: «el disco de centro en O»
106. Con read-arg y prefix: $\spcirc[prefix={el},read-arg={disco}](O)$.

% doc: «circunferencia de centro en O»
107. Con silent false, el valor por defecto: $\spcirc[silent=false](O)$.

% doc: «»
108. Con silent true y sin argumentos: $\spcirc[silent=true]$.

% doc: «»
109. Con silent true y un centro: $\spcirc[silent=true](O)$.

% doc: «»
110. Con silent true y centro y radio: $\spcirc[silent=true](O,3)$.

% doc: «»
111. Con silent true y centro y segmento: $\spcirc[silent=true](O,AB)$.

% doc: «el círculo sub 1 de centro en o subíndice 1 y diámetro medida del lado a b»
112. Con todas las llaves: $\spcirc[read-cmd=mathcat,prefix={el},print-sym=usual,read-sty=circle,diameter=true,radio-sty=mside,concept=true,sub-index={1},read-sub=true,silent=false](O_{1},AB)$.

\section*{Comando spAcirc: argumentos}

% doc: «área de la circunferencia»
113. Sin argumentos: $\spAcirc$.

% doc: «área de la circunferencia con centro en O»
114. Con un centro: $\spAcirc(O)$.

% doc: «área de la circunferencia uno»
115. Con un número: $\spAcirc(1)$.

% doc: «área de la circunferencia con centro en O prima»
116. Con un centro con prima: $\spAcirc(O')$.

% doc: «área de la circunferencia con centro en O sub uno»
117. Con un centro con subíndice: $\spAcirc(O_{1})$.

% doc: «área de la circunferencia con centro en P»
118. Con otra letra como centro: $\spAcirc(P)$.

% doc: «área de la circunferencia dos»
119. Con otro número: $\spAcirc(2)$.

% doc: «área de la circunferencia 12»
120. Con un número de dos cifras: $\spAcirc(12)$.

% doc: «área de la circunferencia n»
121. Con la letra n: $\spAcirc(n)$.

\section*{Comando spAcirc: llaves}

% doc: «área de la circunferencia con centro en O»
122. Con read-cmd school, el valor por defecto: $\spAcirc[read-cmd=school](O)$.

% doc: «área de la circunferencia con centro en o»
123. Con read-cmd mathcat: $\spAcirc[read-cmd=mathcat](O)$.

% doc: «el área de la circunferencia»
124. Con prefix y sin argumentos: $\spAcirc[prefix={el}]$.

% doc: «el área de la circunferencia con centro en O»
125. Con prefix y un centro: $\spAcirc[prefix={el}](O)$.

% doc: «el área de la circunferencia uno»
126. Con prefix y un número: $\spAcirc[prefix={el}](1)$.

% doc: «área de la circunferencia con centro en O»
127. Con read-sty circun, el valor por defecto: $\spAcirc[read-sty=circun](O)$.

% doc: «área del círculo con centro en O»
128. Con read-sty circle: $\spAcirc[read-sty=circle](O)$.

% doc: «área del círculo»
129. Con read-sty circle y sin argumentos: $\spAcirc[read-sty=circle]$.

% doc: «área del círculo uno»
130. Con read-sty circle y un número: $\spAcirc[read-sty=circle](1)$.

% doc: «área de la circunferencia con centro en O»
131. Con concept true, el valor por defecto: $\spAcirc[concept=true](O)$.

% doc: «O»
132. Con concept false y un centro: $\spAcirc[concept=false](O)$.

% doc: «área de la esfera con centro en O»
133. Con read-arg: $\spAcirc[read-arg={esfera}](O)$.

% doc: «esfera con centro en O»
134. Con read-arg y concept false: $\spAcirc[concept=false,read-arg={esfera}](O)$.

% doc: «área de la esfera con centro en O prima»
135. Con read-arg y un centro con prima: $\spAcirc[read-arg={esfera}](O')$.

% doc: «el área de la esfera con centro en O»
136. Con read-arg y prefix: $\spAcirc[prefix={el},read-arg={esfera}](O)$.

% doc: «área del disco con centro en O»
137. Con read-arg y read-sty circle: $\spAcirc[read-sty=circle,read-arg={disco}](O)$.

% doc: «disco con centro en O»
138. Con read-arg y concept false, con otra palabra: $\spAcirc[concept=false,read-arg={disco}](O)$.

% doc: «área de la circunferencia con centro en O»
139. Con print-sym odot, el valor por defecto: $\spAcirc[print-sym=odot](O)$.

% doc: «área de la circunferencia con centro en O»
140. Con print-sym usual: $\spAcirc[print-sym=usual](O)$.

% doc: «área de la circunferencia uno»
141. Con print-sym usual y un número: $\spAcirc[print-sym=usual](1)$.

% doc: «área de la circunferencia»
142. Con print-sym usual y sin argumentos: $\spAcirc[print-sym=usual]$.

% doc: «área de la circunferencia 1 con centro en O»
143. Con sub-index y un centro: $\spAcirc[sub-index={1}](O)$.

% doc: «área de la circunferencia 1 con centro en O prima»
144. Con sub-index y un centro con prima: $\spAcirc[sub-index={1}](O')$.

% doc: «área de la circunferencia sub uno»
145. Con read-sub true y un número: $\spAcirc[read-sub=true](1)$.

% doc: «área de la circunferencia uno»
146. Con read-sub false y un número: $\spAcirc[read-sub=false](1)$.

% doc: «área de la circunferencia sub 1 con centro en O»
147. Con sub-index y read-sub true: $\spAcirc[sub-index={1},read-sub=true](O)$.

% doc: «área de la circunferencia con centro en O»
148. Con silent false, el valor por defecto: $\spAcirc[silent=false](O)$.

% doc: «»
149. Con silent true y sin argumentos: $\spAcirc[silent=true]$.

% doc: «»
150. Con silent true y un centro: $\spAcirc[silent=true](O)$.

% doc: «»
151. Con silent true y un número: $\spAcirc[silent=true](1)$.

% doc: «el área del círculo sub 1 con centro en o»
152. Con todas las llaves: $\spAcirc[read-cmd=mathcat,prefix={el},read-sty=circle,concept=true,print-sym=usual,sub-index={1},read-sub=true,silent=false](O)$.

\section*{Comando spPcirc: argumentos}

% doc: «perímetro de la circunferencia»
153. Sin argumentos: $\spPcirc$.

% doc: «perímetro de la circunferencia con centro en O»
154. Con un centro: $\spPcirc(O)$.

% doc: «perímetro de la circunferencia uno»
155. Con un número: $\spPcirc(1)$.

% doc: «perímetro de la circunferencia con centro en O prima»
156. Con un centro con prima: $\spPcirc(O')$.

% doc: «perímetro de la circunferencia con centro en O sub uno»
157. Con un centro con subíndice: $\spPcirc(O_{1})$.

% doc: «perímetro de la circunferencia con centro en P»
158. Con otra letra como centro: $\spPcirc(P)$.

% doc: «perímetro de la circunferencia dos»
159. Con otro número: $\spPcirc(2)$.

% doc: «perímetro de la circunferencia 12»
160. Con un número de dos cifras: $\spPcirc(12)$.

% doc: «perímetro de la circunferencia n»
161. Con la letra n: $\spPcirc(n)$.

\section*{Comando spPcirc: llaves}

% doc: «perímetro de la circunferencia con centro en O»
162. Con read-cmd school, el valor por defecto: $\spPcirc[read-cmd=school](O)$.

% doc: «perímetro de la circunferencia con centro en o»
163. Con read-cmd mathcat: $\spPcirc[read-cmd=mathcat](O)$.

% doc: «el perímetro de la circunferencia»
164. Con prefix y sin argumentos: $\spPcirc[prefix={el}]$.

% doc: «el perímetro de la circunferencia con centro en O»
165. Con prefix y un centro: $\spPcirc[prefix={el}](O)$.

% doc: «el perímetro de la circunferencia uno»
166. Con prefix y un número: $\spPcirc[prefix={el}](1)$.

% doc: «perímetro de la circunferencia con centro en O»
167. Con read-sty circun, el valor por defecto: $\spPcirc[read-sty=circun](O)$.

% doc: «perímetro del círculo con centro en O»
168. Con read-sty circle: $\spPcirc[read-sty=circle](O)$.

% doc: «perímetro del círculo»
169. Con read-sty circle y sin argumentos: $\spPcirc[read-sty=circle]$.

% doc: «perímetro del círculo uno»
170. Con read-sty circle y un número: $\spPcirc[read-sty=circle](1)$.

% doc: «perímetro de la circunferencia con centro en O»
171. Con concept true, el valor por defecto: $\spPcirc[concept=true](O)$.

% doc: «O»
172. Con concept false y un centro: $\spPcirc[concept=false](O)$.

% doc: «perímetro de la esfera con centro en O»
173. Con read-arg: $\spPcirc[read-arg={esfera}](O)$.

% doc: «esfera con centro en O»
174. Con read-arg y concept false: $\spPcirc[concept=false,read-arg={esfera}](O)$.

% doc: «perímetro de la esfera con centro en O prima»
175. Con read-arg y un centro con prima: $\spPcirc[read-arg={esfera}](O')$.

% doc: «el perímetro de la esfera con centro en O»
176. Con read-arg y prefix: $\spPcirc[prefix={el},read-arg={esfera}](O)$.

% doc: «perímetro del disco con centro en O»
177. Con read-arg y read-sty circle: $\spPcirc[read-sty=circle,read-arg={disco}](O)$.

% doc: «disco con centro en O»
178. Con read-arg y concept false, con otra palabra: $\spPcirc[concept=false,read-arg={disco}](O)$.

% doc: «perímetro de la circunferencia con centro en O»
179. Con print-sym odot, el valor por defecto: $\spPcirc[print-sym=odot](O)$.

% doc: «perímetro de la circunferencia con centro en O»
180. Con print-sym usual: $\spPcirc[print-sym=usual](O)$.

% doc: «perímetro de la circunferencia uno»
181. Con print-sym usual y un número: $\spPcirc[print-sym=usual](1)$.

% doc: «perímetro de la circunferencia»
182. Con print-sym usual y sin argumentos: $\spPcirc[print-sym=usual]$.

% doc: «perímetro de la circunferencia 1 con centro en O»
183. Con sub-index y un centro: $\spPcirc[sub-index={1}](O)$.

% doc: «perímetro de la circunferencia 1 con centro en O prima»
184. Con sub-index y un centro con prima: $\spPcirc[sub-index={1}](O')$.

% doc: «perímetro de la circunferencia sub uno»
185. Con read-sub true y un número: $\spPcirc[read-sub=true](1)$.

% doc: «perímetro de la circunferencia uno»
186. Con read-sub false y un número: $\spPcirc[read-sub=false](1)$.

% doc: «perímetro de la circunferencia sub 1 con centro en O»
187. Con sub-index y read-sub true: $\spPcirc[sub-index={1},read-sub=true](O)$.

% doc: «perímetro de la circunferencia con centro en O»
188. Con silent false, el valor por defecto: $\spPcirc[silent=false](O)$.

% doc: «»
189. Con silent true y sin argumentos: $\spPcirc[silent=true]$.

% doc: «»
190. Con silent true y un centro: $\spPcirc[silent=true](O)$.

% doc: «»
191. Con silent true y un número: $\spPcirc[silent=true](1)$.

% doc: «el perímetro del círculo sub 1 con centro en o»
192. Con todas las llaves: $\spPcirc[read-cmd=mathcat,prefix={el},read-sty=circle,concept=true,print-sym=usual,sub-index={1},read-sub=true,silent=false](O)$.

\section*{Comando spVol}

% doc: «volumen»
193. Sin llaves: $\spVol$.

% doc: «volumen»
194. Con concept true, el valor por defecto: $\spVol[concept=true]$.

% doc: «V»
195. Con concept false: $\spVol[concept=false]$.

% doc: «V»
196. Con concept false y read-cmd school: $\spVol[concept=false,read-cmd=school]$.

% doc: «guión v»
197. Con concept false y read-cmd mathcat: $\spVol[concept=false,read-cmd=mathcat]$.

% doc: «volumen 1»
198. Con sub-index: $\spVol[sub-index={1}]$.

% doc: «volumen 12»
199. Con sub-index de dos cifras: $\spVol[sub-index={12}]$.

% doc: «volumen sub 1»
200. Con sub-index y read-sub true: $\spVol[sub-index={1},read-sub=true]$.

% doc: «volumen 1»
201. Con sub-index y read-sub false: $\spVol[sub-index={1},read-sub=false]$.

% doc: «V 1»
202. Con sub-index y concept false: $\spVol[sub-index={1},concept=false]$.

% doc: «volumen medida exacta»
203. Con read-arg: $\spVol[read-arg={medida exacta}]$.

% doc: «medida exacta»
204. Con read-arg y concept false: $\spVol[concept=false,read-arg={medida exacta}]$.

% doc: «volumen medida exacta»
205. Con read-arg y sub-index: $\spVol[read-arg={medida exacta},sub-index={1},read-sub=true]$.

% doc: «volumen»
206. Con silent false, el valor por defecto: $\spVol[silent=false]$.

% doc: «»
207. Con silent true: $\spVol[silent=true]$.

% doc: «volumen subíndice 1»
208. Con todas las llaves: $\spVol[concept=true,read-cmd=mathcat,sub-index={1},read-sub=true,silent=false]$.

\section*{En fórmulas}

% doc: «diámetro es igual a 2 punto radio»
209. Diámetro y radio: $\spdiam = 2 \cdot \sprad$.

% doc: «área de la circunferencia con centro en O perímetro de la circunferencia con centro en O»
210. Área y perímetro de la misma circunferencia: $\spAcirc(O) \quad \spPcirc(O)$.

% doc: «circunferencia de centro en O y radio 3 flecha doble derecha radio es igual a 3»
211. Una circunferencia y su radio: $\spcirc(O,3) \Rightarrow \sprad = 3$.

% doc: «perímetro de la circunferencia con centro en O es igual a 2 punto pi punto radio»
212. Perímetro en función del radio: $\spPcirc(O) = 2 \cdot \pi \cdot \sprad$.

% doc: «área de la circunferencia con centro en O es igual a pi punto radio»
213. En fórmula destacada:
\[ \spAcirc(O) = \pi \cdot \sprad \].

% doc: «volumen es igual a 4 tercios punto pi punto radio punto radio punto radio»
214. Volumen de una esfera sin potencias: $\spVol = \frac{4}{3} \cdot \pi \cdot \sprad \cdot \sprad \cdot \sprad$.

\section*{Casos con potencias (al final, por precaución)}

% doc: «área de la circunferencia con centro en O es igual a pi radio cuadrado»
215. Área con exponente: $\spAcirc(O) = \pi \sprad^{2}$.

% doc: «volumen es igual a 4 tercios pi radio al cubo»
216. Volumen con exponente: $\spVol = \frac{4}{3} \pi \sprad^{3}$.

% doc: «diámetro cuadrado es igual a 4 radio cuadrado»
217. Diámetro al cuadrado: $\spdiam^{2} = 4 \sprad^{2}$.

\end{document}
